\documentclass[10pt, aspectratio=169]{beamer}

% Konfiguracja języka i kodowania
\usepackage[utf8]{inputenc}
\usepackage[T1]{fontenc}
\usepackage[polish]{babel}

\usetheme{Madrid}
\usecolortheme{whale}


\usepackage{listings}
\usepackage{xcolor}
\usepackage{tikz}
\usetikzlibrary{automata, positioning, arrows, arrows.meta}

\definecolor{bluekeywords}{rgb}{0,0,1}
\definecolor{greencomments}{rgb}{0,0.5,0}
\definecolor{redstrings}{rgb}{0.64,0.08,0.08}
\definecolor{xmlcomments}{rgb}{0.5,0.5,0.5}

\setbeamerfont{title}{size=\large}
\title{Naprawianie automatów skończonych z brakującymi stanami i krawędziami}
\subtitle{Repairing Finite Automata with Missing Components}
\author{Julia Cygan, Patryk Flama}
\institute{II UWr}
\date{\today}

\begin{document}

% Slajd tytułowy
\begin{frame}
  \titlepage
\end{frame}

% % Spis treści
% \begin{frame}{Spis treści}
%   \tableofcontents
% \end{frame}

\section{Definicja problemu}
\begin{frame}{Definicja problemu}

\only<1>{
  \begin{columns}[c]
    \column{0.48\textwidth}
    \begin{block}{DFA}
      Deterministyczny automat skończony to krotka:
      $$(Q, \Sigma, \delta, q_\lambda, \mathbb{F_A}, \mathbb{F_R})$$
      gdzie:
      \begin{itemize}
        \small
        \item $Q$ - zbiór stanów
        \item $\Sigma$ - alfabet
        \item $\delta: Q \times \Sigma \to Q$ - przejścia
        \item $q_\lambda$ - stan początkowy
        \item $\mathbb{F_A}$ - stany akceptujące
        \item $\mathbb{F_R}$ - stany odrzucające
      \end{itemize}
    \end{block}
    
    \column{0.48\textwidth}
    \vfill
    \begin{center}
      \begin{tikzpicture}[
        node distance=2cm,
        initial text=,
        every state/.style={minimum size=0.8cm}
      ]
        \node[state] (q0) {$q_\lambda$};
        \node[state, accepting, right of=q0] (q1) {$q_1$};
        \node[state, right of=q1] (q2) {$q_2$};
        
        \path[->]
          (q0) edge[above] node{$a$} (q1)
          (q1) edge[above] node{$a,b$} (q2)
          (q2) edge[loop right] node{$a,b$} (q2)
          (q0) edge[bend right, below] node{$b$} (q2);
      \end{tikzpicture}
    \end{center}
    \vfill
  \end{columns}
}

\only<2>{
  \begin{columns}[c]
    \column{0.52\textwidth}
    \begin{block}{Z nieokreśloną liczbą stanów}
      \small
      \textbf{Wejście:} Częściowy DFA $A$ z brakującymi krawędziami i stanami, zbiory $S^+$, $S^-$.
      
      \textbf{Wyjście:} Czy można uzupełnić przejścia, klasyfikacje stanów i dodać stany, aby automat był poprawny?
    \end{block}
    
    \vspace{0.3cm}
    
    \begin{block}{Z określoną liczbą stanów}
      \small
      \textbf{Wejście:} Częściowy DFA $A$ z brakującymi krawędziami, zbiory $S^+$, $S^-$.
      
      \textbf{Wyjście:} Czy można uzupełnić przejścia i klasyfikacje stanów, aby automat był poprawny?
      
      (Bez dodawania nowych stanów)
    \end{block}
    
    \column{0.48\textwidth}
    \vfill
    \begin{center}
      \begin{tikzpicture}[
        ->,
        >=Stealth,
        node distance=2.2cm,
        initial text=,
        every state/.style={circle, draw, minimum size=0.9cm}
      ]
        \node[state] (q0) at (0,0) {$q_\lambda$};
        \node[state] (q1) at (2.2,0) {$q_1$};
        \node[state, accepting] (q2) at (4.4,0) {$q_2$};
        
        \coordinate (undef1) at (0,-1.5);
        \coordinate (undef2) at (2.2,1.5);
        
        \draw[->,thick] (q0) -- (q1) node[midway,above] {$a$};
        \draw[->,thick] (q1) -- (q2) node[midway,above] {$a,b$};
        \draw[->,thick] (q2) edge[loop right] node{$a,b$} (q2);
        \draw[->,dashed,thick] (q0) -- (undef1) node[midway,right] {$b$};
        \draw[->,dashed,thick] (q1) -- (undef2) node[midway,right] {$a$};
      \end{tikzpicture}
    \end{center}
    \vfill
  \end{columns}
}

\end{frame}

% TODO: fajny pomysł na slajd, ale content słaby
\begin{frame}{Dlaczego ten problem jest ważny?}
  \begin{block}{Znany trudny problem}
    Problem \textit{DFA-consistency} (najmniejszy zgodny automat):
    \begin{itemize}
      \item \textbf{Wejście:} liczba $n$, zbiory $S^+$, $S^-$
      \item \textbf{Wyjście:} Czy istnieje DFA z $\leq n$ stanami akceptujący $S^+$ i odrzucający $S^-$?
      \item \textbf{Znany wynik:} $\mathbb{NP}$-zupełny
    \end{itemize}
  \end{block}
  
  \vspace{0.3cm}
  
  \begin{block}{Nasz problem}
    \begin{itemize}
      \item Generalnie co najmniej tak trudny co DFA-consistency
      \item Można pokazać redukcję z DFA-consistency
      \item Silnie motywuje teoretyczną analizę naszego problemu
    \end{itemize}
  \end{block}
\end{frame}

% TODO: też jakiś refactor - za dużo słów + jakiś wstęp po co rozważać złożoność
%* dodatkowo fajnie powiedzieć że:
%* - problem z nieograniczoną liczbą stanów jest trudniejszy
%* - problem z nieograniczoną liczbą stanów potrafimy zredukować liniowo do problemu z ograniczoną liczbą stanów (w zależności od rozmiaru próbek wejściowych)
%* - dlatego wystarczy nam rozważać problem z ograniczoną liczbą stanów
\section{Złożoność obliczeniowa}
\begin{frame}{NP-zupełność}
  \begin{block}{Cel: wykazać, że problem jest NP-zupełny}
    Trzeba pokazać dwie rzeczy:
    \begin{enumerate}
      \item Problem jest w $\mathbb{NP}$ (certyfikat istnieje)
      \item Problem jest $\mathbb{NP}$-trudny (nie ma algorytmu wielomianowego)
    \end{enumerate}
  \end{block}
  
  \vspace{0.1cm}
  
  \begin{block}{Problem jest w $\mathbb{NP}$}
    Dla danego rozwiązania (uzupełnienia automatu) weryfikacja zajmuje czas wielomianowy:
    \begin{itemize}
      \small
      \item Determinizm automatu: $O(|Q| \cdot |\Sigma|)$
      \item Sprawdzenie próbek: $O((|S^+| + |S^-|) \cdot M)$ gdzie $M$ to max długość
    \end{itemize}
  \end{block}
  
  \vspace{0.1cm}
  
  \begin{block}{Problem jest $\mathbb{NP}$-trudny}
    \textbf{Redukcja z DFA-consistency:}
    
    Dla instancji $(n, S^+, S^-)$ tworzymy częściowy automat z $n$ stanami, wszystkimi przejściami oraz wartościowaniami stanów niezdefiniowanymi. Jeśli można go naprawić, to rozwiązuje DFA-consistency.
  \end{block}
\end{frame}

\end{document}
